
% language=us runpath=t:/texmf/doc/context/documents

% This is a very old style, only adapted to the new xml handler. Nowdays articles
% are collected in documents. We started with \TEX\ files, later moved to \XML\
% with \MKII\ processing (using a streaming parser and save/reuse data model).
%
% In order to also discuss the \LMTX\ manuals we reorganized the \XML\ file a bit
% the style is mostly unchanged. We did change the font to pennstander.
%
% We assume that this manual is processed here: texmf/doc/context/documents
%
% A local helper: t:\manuals\mkii\external\tempcopy.cmd

% we can reuse some covers:
%
% (m)pstopdf
% (m)texsync

\environment show-env

% \TransparencyHack

\setuplayout
  [topspace=30pt,
   height=fit,
   backspace=30pt,
   width=300pt,
   rightmargin=220pt,
   margindistance=30pt,
   header=0pt,
   footer=0pt,
   bottomdistance=15pt,
   bottom=20pt]

\defineoverlay
  [MainBackground]
  [\useMPgraphic{MainGraphic}]

\setupbackgrounds
  [page]
  [background={content,MainBackground}]

\setupcolors
  [textcolor=white]

\defineoverlay[content][\overlaybutton{content}]

\startuseMPgraphic{MainGraphic}

    % todo: adapt the shades, use the luametafun variant:

    StartPage ;

        path p ;

        fill Page withcolor "shade-2" ;

        p := llcorner Field[Text][Bottom] --
             lrcorner Field[Text][Bottom] --
             urcorner Field[Text][Text]   --
             ulcorner Field[Text][Text]   -- cycle ;
        p := p enlarged .5cm randomized .5cm ;

        fill p
            withshademethod "circular"
            withshadecolors ("shade-1", "shade-2") ;

        def bottom_menu_button (expr nn, rr, pp, xx, yy, ww, hh, dd) =
            if (pp > 0) and (rr > 0) :
                if     nn = 1 :
                     p := fullsquare xysized (ww/2,hh) ;
                     p := p shifted (xx+ww/4,yy+hh/2) ;
                elseif nn = 3 :
                     p := fulltriangle xysized (ww,hh) ;
                     p := p shifted (xx+ww/2,yy+hh/2) ;
                elseif nn = 2 :
                     p := fulltriangle rotated 180 xysized (ww,hh) ;
                     p := p shifted (xx+ww/2,yy+hh/2) ;
                else :
                     p := origin -- cycle ;
                fi ;
                 p := p randomizedcontrols (hh/4) ;
                 fill p
                    withshademethod "circular"
                    withshaderadius (0,ww)
                    withshadecolors ("shade-1", "shade-2") ;
            fi ;
        enddef ;

        \MPmenubuttons{bottom}

    StopPage ;

\stopuseMPgraphic

\setupinteractionmenu
  [bottom]
  [state=start,
   frame=off,
   left=\hfill,
   middle=\hskip.5cm,
   width=2\bottomheight,
   position=yes]

\startinteractionmenu[bottom]
    \startbut [firstpage]    \stopbut
    \hfilll
    \startbut [previouspage] \stopbut
    \startbut [nextpage]     \stopbut
\stopinteractionmenu

\newbox\samplebox

\setupexternalfigures[compact=]

\definelist
  [manual]
  [alternative=command]

\starttexdefinition SomePlaceholder #1
    \writestatus{NOT FOUND}{\ThePaper.pdf}
    %
    \edef\StyleWidth{\the\dimexpr(\textwidth-60pt)/6\relax}% 7
    %
    \framed [
        offset=overlay,
        frame=off,
        background=color,
        backgroundcolor=shade-2,
        backgroundoffset=2.5pt,
        width=\StyleWidth,
        height=29.7\dimexpr\StyleWidth/21\relax
    ] { \smallinfofont
        \cldcontext{file.basename("#1")}
    }
\stoptexdefinition

\globaldisablemode[continuation]

\startxmlsetups xml:overview

    \edef\StyleWidth{\the\dimexpr(\textwidth-60pt)/6\relax}% 7
    \edef\ThePaper  {\xmlflush{#1}}% weird. why no xmltext

    \dontleavehmode
    \doifelsefile{\ThePaper.pdf} {
        \button [
            offset=overlay,
            frame=off,
            background=color,
            backgroundcolor=shade-2,
            backgroundoffset=2.5pt
        ] {
            \externalfigure [\ThePaper.pdf] [
                background=color,
                backgroundcolor=white,
                page=1,
                width=\StyleWidth,
                height=29.7\dimexpr\StyleWidth/21\relax
            ]
        } [
            \ThePaper
        ]

    } {
        \SomePlaceholder{\ThePaper}
    }
    \hskip10pt plus 10pt\relax

    \allowbreak

\stopxmlsetups

\startxmlsetups xml:margintext
    \vbox to \textheight {
        \hsize\rightmarginwidth
        \setbox\samplebox\vbox {
            \setupwhitespace[big]
            \setupalign[flushleft,mathbookpasses]
            \bgroup
                \bfd
                \setupinterlinespace
                \xmlfirst{#1}{/title}
                \par
            \egroup
            \blank
            \bgroup
                \bf
                \setupinterlinespace
                \xmlfirst{#1}{/summary}
                \par
                \doifmode {continuation} {
                    (continued)
                }
            \egroup
        }
        \button
            [frame=off,offset=overlay,style=,color=]
            {\box\samplebox}
            [forward]
        \vfill
    }

    \globalenablemode[continuation]

\stopxmlsetups

\startxmlsetups xml:documents

    % common

    \globaldisablemode[continuation]

    \definecolor[dark][s=.25]

    \setuptexttexts[margin][][\pagereference[all]\xmlfilter{#1}{/introduction[@version=='all']/command(xml:margintext)}]

    \start
        \start
            \setupbutton
              [color=white,
               contrastcolor=white,
               width=1tw,
               background=color,
               frame=off,
               backgroundcolor=dark,
               style=\bfc]

            \button{common manuals}[common] \blank
            \button{\MKXL\ manuals}[mkxl]   \blank
            \button{\MKIV\ manuals}[mkiv]   \blank
            \button{\MKII\ manuals}[mkii]
        \stop
        \page
    \stop

    \globaldisablemode[continuation]

    \setuptexttexts[margin][][\pagereference[common]\xmlfilter{#1}{/introduction[@version=='common']/command(xml:margintext)}]

    \start
        \setupalign[flushleft,mathbookpasses]
        \lineskip10pt
        \xmlfilter{#1}{/collection[@version=='common' and @classification!='old']/resources/manual[match()==1]/paper/command(xml:overview)}%
        \page
    \stop

    % mkxl

    \globaldisablemode[continuation]

    \setuptexttexts[margin][][\pagereference[mkxl]\xmlfilter{#1}{/introduction[@version=='mkxl']/command(xml:margintext)}]

    \start
        \setupalign[flushleft,mathbookpasses]
        \lineskip10pt
        \xmlfilter{#1}{/collection[@version=='mkxl' and @classification!='old']/resources/manual[match()==1]/paper/command(xml:overview)}%
        \page
    \stop

    % mkiv

    \globaldisablemode[continuation]

    \setuptexttexts[margin][][\pagereference[mkiv]\xmlfilter{#1}{/introduction[@version=='mkiv']/command(xml:margintext)}]

    \start
        \setupalign[flushleft,mathbookpasses]
        \lineskip10pt
        \xmlfilter{#1}{/collection[@version=='mkiv' and @classification!='old']/resources/manual[match()==1]/paper/command(xml:overview)}%
        \page
    \stop

    % mkii

    \globaldisablemode[continuation]

    \setuptexttexts[margin][][\pagereference[mkii]\xmlfilter{#1}{/introduction[@version=='mkii']/command(xml:margintext)}]

    \start
        \setupalign[flushleft,mathbookpasses]
        \lineskip10pt
        \xmlfilter{#1}{/collection[@version=='mkii' and @classification!='old']/resources/manual[match()==1]/paper/command(xml:overview)}%
        \page
    \stop

    % obsolete

    \globaldisablemode[continuation]

    \setuptexttexts[margin][][\xmlfilter{#1}{/introduction[@classification=='old']/command(xml:margintext)}]

    \start
        \setupalign[flushleft,mathbookpasses]
        \lineskip10pt
        \xmlfilter{#1}{/collection[@classification=='old']/resources/manual[match()==1]/paper/command(xml:overview)}%
        \page
    \stop

    \setuptexttexts[margin][][]

    \globaldisablemode[continuation]

    % everything

    \xmlall{#1}{/collection}

\stopxmlsetups

\startxmlsetups xml:sample

    \edef\ThePaper {\xmlflush{#1}}
    \edef\TheColor {\xmlatt{#1}{backgroundcolor}}
    \edef\TheOffset{\xmlatt{#1}{offset}}

    \doifelse {\TheOffset} {yes} {
        \def\TheOffset{.25ex}
    } {
        \def\TheOffset{0pt}
    }

    \doifelsefile{\ThePaper.pdf} {
        \gotobox {
            \doifelsenothing {\TheColor}
                {\externalfigure[\ThePaper.pdf][page=1,width=\StampWidth]}
                {\externalfigure[\ThePaper.pdf][page=1,width=\StampWidth,background=color,backgroundcolor=\TheColor,backgroundoffset=\TheOffset]}
        }
        [file(\ThePaper)]
    } {
        \SomePlaceholder{\ThePaper}
    }

    \hskip\StampDelta

\stopxmlsetups

\startxmlsetups xml:collection

    \edef\FirstRow {\xmlcount{#1}{resources/manual/paper}}
    \edef\SecondRow{\xmlcount{#1}{resources/manual/screen}}

    \setbox\samplebox=\vbox to \textheight {

        \hsize\rightmarginwidth

        \vbox \bgroup
            \raggedright
            \bgroup
                \bfd
                \setupinterlinespace
                \xmlfirst{#1}{description/theme}
                \par
            \egroup
            \blank
            \bf
            \xmlfirst{#1}{description/summary}
            \blank

                                \def\StampWidth{45pt} \def\StampDelta {10pt}
            \ifnum  \FirstRow>4 \def\StampWidth{35pt} \def\StampDelta {10pt} \fi
            \ifnum \SecondRow>4 \def\StampWidth{35pt} \def\StampDelta {10pt} \fi
            \ifnum  \FirstRow>6 \def\StampWidth{25pt} \def\StampDelta {10pt} \fi
            \ifnum \SecondRow>6 \def\StampWidth{25pt} \def\StampDelta {10pt} \fi

            \ifnum  \FirstRow>1 \vskip10pt \hbox to \hsize{\hss\xmlfilter{#1}{resources/manual[position()>1]/paper/command(xml:sample)}\unskip} \fi
            \ifnum \SecondRow>1 \vskip10pt \hbox to \hsize{\hss\xmlfilter{#1}{resources/manual/screen/command(xml:sample)}\unskip} \fi
        \egroup

        \vfill

        \edef\TheScreen{\xmlfilter{#1}{resources/manual/screen[position()=1]/text()}}
        \edef\TheColor {\xmlfilter{#1}{resources/manual/screen/attribute(backgroundcolor)}}

        \doifsomething {\TheScreen} {
            \hfill
            \gotobox {
                \doifelsenothing {\TheColor} {
                    \externalfigure[\TheScreen.pdf][page=1,width=90pt]
                } {
                    \externalfigure[\TheScreen.pdf][page=1,width=90pt,background=color,backgroundcolor=\TheColor]
                }
            }
            [\TheScreen::]
        }

        \vskip-\bottomheight
        \vskip-\bottomdistance
    }

    \setuptexttexts[margin][][\box\samplebox]

    \startstandardmakeup

        \edef\ThePaper{\xmlfilter{#1}{resources/manual/paper[position()=1]/text()}}
        \edef\TheColor{\xmlfilter{#1}{resources/manual/paper/attribute(backgroundcolor)}}

        \centerbox {
            \gotobox {
                \doifelsenothing {\TheColor} {
                    \externalfigure[\ThePaper.pdf][page=1,factor=max]
                } {
                    \externalfigure[\ThePaper.pdf][page=1,factor=max,background=color,backgroundcolor=\TheColor]
                }
            }
            [\ThePaper::]
        }

        \pagereference[\ThePaper]

    \stopstandardmakeup

\stopxmlsetups

\startxmlsetups xml:definitions
    \xmlsetsetup{#1}{collection|documents|introduction}{xml:*}
\stopxmlsetups

\xmlregistersetup{xml:definitions}

\starttext
    \xmlprocess{main}{show-man.xml}{}
\stoptext
